\chapter{Implementation and Results}
This is the 1st work.
\section{Performance metrics}
System performance was evaluated using widely accepted classification and
explainability metrics suitable for healthcare-oriented artificial intelligence
systems. These metrics collectively assess predictive accuracy, robustness of
visual preprocessing, and transparency of decision-making.
The system performance was evaluated using standard classification and
explainability metrics commonly adopted in healthcare-oriented AI
systems. Table \ref{tab:performance_metrics} summarises the evaluation
criteria applied in this study.
\begin{table}[h]
\centering
\begin{tabular}{|l|p{9cm}|}
\hline
\textbf{Metric} & \textbf{Description} \\ \hline
F1-Score (Macro) & Measures classification balance across all ASD behavioural classes \\ \hline
ROC-AUC & Evaluates the discriminative capability of the classifier \\ \hline
SSIM & Measures structural quality of preprocessed visual frames \\ \hline
Explainability Index (EI) & Quantifies clarity and usefulness of generated explanations \\ \hline
Faithfulness Score & Measures alignment between explanations and model decision process \\ \hline
\end{tabular}
\caption{Performance Metrics Used for Evaluation}
\label{tab:performance_metrics}
\end{table}

\section{test cases}
System performance was evaluated using widely accepted classification and
explainability metrics suitable for healthcare-oriented artificial intelligence
systems. These metrics collectively assess predictive accuracy, robustness of
visual preprocessing, and transparency of decision-making.
The system performance was evaluated using standard classification and
explainability metrics commonly adopted in healthcare-oriented AI
systems.summarises the evaluation
criteria applied in this study.


